\chapter{Informe de Gestión del Presidente (Oct 2023 -- Oct 2025)}

Este período tuvo tres etapas claras. En 2023 nos enfocamos en estabilizar lo urgente y conocer a fondo el estado del sistema. En 2024 pasamos de reaccionar a mejorar: abrimos cajas y válvulas para sectorizar, coordinamos con instituciones y reforzamos el control de la calidad del agua. En 2025 dimos el paso que faltaba: ordenar mejor la cloración y, sobre todo, levantar un estudio técnico para contar con planos y datos que permitan decidir con fundamento

\section{Año 2023}

Arrancamos la gestión atendiendo lo urgente: asegurar operación diaria, conocer el estado real del sistema y ordenar lo básico en administración. Formalizamos la contratación del operador para tener presencia técnica estable, pusimos al día la cloración y recorrimos captación y planta para entender dónde estaban los cuellos de botella. También hicimos mejoras puntuales en la red para que los arreglos sean más rápidos y seguros

\subsection*{Qué hicimos y por qué}
\begin{itemize}[noitemsep, topsep=0pt]
  \item \textbf{Contrato del operador} con jornada de 40 horas, afiliación al IESS y responsabilidades claras, para que el sistema tenga seguimiento diario y trazabilidad
  \item \textbf{Mantenimiento de la cloradora} con reposición y calibración, porque sin dosificación estable no hay agua segura
  \item \textbf{Inspección a planta y captación} para revisar estructuras, caudales y pérdidas. En la captación observamos caudal reducido por estiaje y tramos prácticamente secos; planificamos limpieza de rejillas, control de sólidos y mediciones de caudal para ajustar la cloración
  \item \textbf{Cambio de tapa de caja de revisión} en Barrio Progreso, mejorando seguridad y acceso para maniobras y reparaciones
  \item \textbf{Mejora de filtración} con cambio de medios filtrantes en filtros lento y grueso, combinando minga y contratación para recuperar continuidad y calidad
  \item \textbf{Apoyo contable y de inventarios} mediante proyecto de vinculación, para ordenar registros físicos y digitales
\end{itemize}

\subsection*{Cronología 2023}
\begin{table}[H]
\centering
\caption{Cronología de actividades 2023}
\label{tab:cronologia_2023}
\begin{singlespace}
\begin{tabular}{p{3cm} p{11cm}}
\toprule
\textbf{Fecha} & \textbf{Actividad} \\
\midrule
04 sep 2023 & Cambio de tapa de caja de revisión en Barrio Progreso; priorizamos seguridad y tiempos de intervención \\
22 sep 2023 & Inspección en planta de tratamiento; se definieron mejoras en filtros y cajas de entrada y salida \\
02 oct 2023 & Contrato del operador con horario y funciones definidos para asegurar continuidad operativa \\
03 oct 2023 & Mantenimiento y calibración de la máquina clorificadora para estabilizar la dosificación \\
20 nov 2023 & Inspección a la captación; se evidenció estiaje, ingreso de sólidos y necesidad de proteger la toma \\
20 nov -- 16 dic 2023 & Mejora de filtración con cambio de medio filtrante, combinando minga y contratación puntual \\
dic 2023 & Orden administrativo y soporte contable mediante proyecto de vinculación \\
\bottomrule
\end{tabular}
\end{singlespace}
\end{table}

\section{Año 2024}

En 2024 pasamos de la emergencia a la mejora. Trabajamos en la red para reducir pérdidas y tiempos de reparación, coordinamos con instituciones para estudios y materiales, y reforzamos el control de calidad del agua. También abrimos espacios con la comunidad para explicar lo que hacemos y por qué es importante cuidar el sistema

\subsection*{Intervenciones y gestiones clave}
\begin{itemize}[noitemsep, topsep=0pt]
  \item \textbf{Cajas de revisión y válvulas} en Tres Esquinas, Centro y Progreso, para sectorizar mejor y cerrar más rápido cuando hay reparaciones
  \item \textbf{Coordinación con el GAD Municipal} para la ficha técnica de ampliación de matriz; asignación de personal técnico para el levantamiento
  \item \textbf{Visitas técnicas} a red, captación y planta; constatamos cloración insuficiente y se recomendó migrar a hipoclorito de calcio, elaborar mapa con coordenadas y medir caudales para ajustar dosis
  \item \textbf{Control y orden} en la red con sanción a una conexión clandestina y reparación puntual en Barrio San Juan
  \item \textbf{Conservación de microcuencas} con minga de siembra de plantas en riberas del río Cachi, articulando con EMAIPC y comunidad
  \item \textbf{Gestión de espacios} con renovación de comodato de oficinas por cinco años, para asegurar continuidad administrativa
\end{itemize}

\subsection*{Cronología 2024}
\begin{table}[H]
\centering
\caption{Cronología de actividades 2024}
\label{tab:cronologia_2024}
\begin{singlespace}
\begin{tabular}{p{3cm} p{11cm}}
\toprule
\textbf{Mes} & \textbf{Actividad} \\
\midrule
Enero & Apoyo a obra vial en La Campanera coordinando materiales y mano de obra para proteger instalaciones de agua \\
Febrero & Revisión de desfogue en Barrio La Loma y verificación de válvulas y accesorios \\
Marzo & Inicio de construcción de cajas de revisión y cambio de válvulas en sectores clave para sectorización y disminución de pérdidas \\
Abril & Continuación de obras y gestión ante GAD Municipal para ficha técnica de ampliación de matriz \\
Mayo & Coordinación con GAD Parroquial por alcantarillado pluvial; visita a PTAR La Carmela; registro de nuevo derecho de agua y actualización de padrón \\
Junio & Participación cívica para fortalecer relación con usuarios y actores locales \\
Julio & Recorridos técnicos a red, captación y planta; se recomienda migrar a hipoclorito de calcio, hacer mapa con coordenadas y medir caudales para ajustar dosificación \\
Agosto & Actividades cívicas cantonales y relación comunitaria \\
Septiembre & Inicio de la Escuelita del Agua; socialización de obras de asfaltado y revegetación con EMAIPC \\
Octubre & Inspección del proyecto de asfaltado; sanción por conexión clandestina; reparación de red en Barrio San Juan \\
Noviembre & Solicitud de materiales al GAD Municipal; intercambio técnico con otra Junta; mediciones de caudal en sequía; avance para certificado ambiental \\
Diciembre & Minga general con siembra de plantas en riberas del río Cachi; renovación del comodato de oficinas por cinco años \\
\bottomrule
\end{tabular}
\end{singlespace}
\end{table}

\section{Año 2025}

En 2025 dimos orden a la caseta de cloración y, sobre todo, avanzamos con el estudio técnico con el Ing.\ Remigio Ojeda. Esto nos permitirá por primera vez contar con planos y con el \emph{esqueleto} de la red para decidir con fundamento. Al mismo tiempo atendimos derrumbes y roturas, principalmente en La Campanera y en el barrio San Juan, lo que mostró la necesidad de proteger mejor la infraestructura y de seguir sectorizando con criterio. Además, iniciamos la reforma del estatuto y del reglamento para actualizar reglas internas y responsabilidades

\subsection*{Intervenciones y gestiones clave}
\begin{itemize}[noitemsep, topsep=0pt]
  \item Mejora de la caseta de cloración con dos mesones para ordenar tanques y equipos, facilitando mantenimiento y seguridad
  \item Limpieza de captación y desarenador para reducir la entrada de sedimentos a la planta
  \item Atención de derrumbes y roturas en La Campanera y en el barrio San Juan, priorizando continuidad del servicio y protección de la línea
  \item Estudio técnico con el Ing.\ Remigio Ojeda: recorridos, toma de muestras, pruebas de proceso y socializaciones para levantar información completa y proponer mejoras con base técnica
  \item Trabajo comunitario y coordinación con el GAD Parroquial y el GAD Municipal para presupuesto, apoyo a usuarios y cierre de capacitaciones
  \item Inicio de la reforma del estatuto y del reglamento con asesoría legal, para fortalecer la gestión y el control interno
\end{itemize}

\subsection*{Cronología 2025}
\begin{table}[H]
\centering
\caption{Cronología de actividades 2025}
\label{tab:cronologia_2025}
\begin{singlespace}
\begin{tabular}{p{3cm} p{11cm}}
\toprule
\textbf{Mes} & \textbf{Actividad} \\
\midrule
Enero & Mejora en la caseta de cloración: construcción de dos mesones para reubicar tanques y ordenar equipos \\
Febrero & Limpieza de la captación y del desarenador; cierre de actividades de la Escuelita del Agua con participación comunitaria; deslizamientos cercanos a la captación y afectaciones en La Campanera \\
Marzo & Asamblea para el informe de actividades 2024 (fecha por confirmar, primera quincena); firma del contrato con el Ing.\ Remigio Ojeda para el estudio de la red \\
Abril & Nuevos derrumbes en La Campanera; toma de muestras de agua para el estudio; entrega de pollos a usuarios en proyecto del GAD Municipal y Elecaustro \\
Mayo & Recorrido con el Ing.\ Ojeda para evaluar una posible nueva captación en el sector Doña Justa \\
Junio & Asamblea para explicar el estudio contratado; falla en el barrio San Juan y reparación de la línea; participación en el desfile cívico de Jerusalén \\
Julio & Afectaciones en La Campanera por inestabilidad del talud; pruebas en planta para evaluar mejoras del proceso; primeras socializaciones del estudio con la comunidad \\
Agosto & Reunión con el GAD Parroquial por presupuesto participativo; entrega del proyecto técnico por parte del Ing.\ Ojeda \\
Septiembre & Solicitud de presupuesto adicional ante el GAD de Biblián para implementar mejoras; trabajos de pintura y orden de instalaciones \\
Octubre & Falla en el barrio San Juan con corrección de la línea; nuevas afectaciones en La Campanera por lluvias y movimiento de suelo. Entrega de los nuevos estatutos y reglamentos para trámite ante la autoridad competente \\
\bottomrule
\end{tabular}
\end{singlespace}
\end{table}

\subsection*{Cierre del período y próximos pasos}
Este período nos deja un aprendizaje claro: no basta con reparar, hay que anticiparse con información. El estudio técnico y los planos serán la base para un plan de inversiones por fases, con sectorización, protección de taludes y renovación de tramos críticos. En planta, las pruebas y el orden en la caseta de cloración facilitan operaciones más seguras y estables. La coordinación con el GAD Parroquial, el GAD Municipal y programas aliados permitirá gestionar recursos para las mejoras priorizadas. Los detalles económicos se presentan en el Capítulo 2 para mantener separadas la narrativa de gestión y las cifras
